Predicting the semantic orientation of adjectives \cite{hatzivassiloglou1997predicting}
Learning subjective adjectives from corpora \cite{wiebe2000learning}
Thumbs up or thumbs down?: semantic orientation applied to unsupervised classification of reviews \cite{turney2002thumbs}
Trained a machine learning classifier \cite{pang2002thumbs}

\cite{turney2003measuring}
\cite{taboada2004analyzing}

Construct sentiment dictionary from lexical graph \cite{kim2004determining}
Construct sentiment dictionary from lexical graph \cite{hu2004mining}

Manually constructed polarity lexicon \cite{das2007yahoo}

Building Lexicon for Sentiment Analysis from Massive Collection of HTML Documents. \cite{kaji2007building}

Construct sentiment dictionary from lexical graph \cite{blair2008building}

\cite{bestgen2008building}

Contrastive summarization: an experiment with consumer reviews\cite{lerman2009contrastive}

Active learning for text classification \cite{esuli2009active}
Semi-supervised learning of lexicon \cite{rao2009semi}

\cite{velikovich2010viability}
The viability of web-derived polarity lexicons
Used n-grams from 4 billion web pages, and pruned a network with various heuristics contructed from context vectors (window size 6).
Seeds words for positive and negative sentiment, and scores determined by shortest paths to the seed set, a ``graph propogation algorithm''.
Compares favorably to manually constructed lexicon from \cite{wilson2005recognizing} and lexicon from WordNet in \cite{blair2008building}.

\cite{bestgen2012checking}

\cite{carenini2013multi}
Multi-document summarization of evaluative text

\cite{tang2014building}

Other neural nets:
\cite{kim2014convolutional}

\cite{recchia2015reproducing}

\cite{amir2016expanding}
Expanding Subjective Lexicons for Social Media Mining with Embedding Subspaces
Their approach to lexicon expansion ``consists of training models to predict the labels of pre-existing lexicons, leveraging unsupervised word embeddings as features'' \cite{amir2016expanding}.
Correlations between their method and existing continuous datasets had a maximum of 0.68, an improvement over support vector regression.
The resulting lexicon out-performed the methods compared against in Tweet classification, although not all methods were compared.

\cite{hamilton2016inducing}
Inducing Domain-Specific Sentiment Lexicons from Unlabeled Corpora
Utilizes the approach set out in \cite{velikovich2010viability} to generate corpus specific word embeddings using SVD and propograte sentiment labels on inferred $k$-NN network.
Measures uncertainty in predicted labels with bootstrapping procedure that holds out fractions of seed set (on the order of 10 words, holding out 2).
Claims to measure performance with correlations to existing dataset of \cite{warriner2013norms}, but not found in results.

\cite{williams2016boundary}
Boundary-based MWE segmentation with text partitioning
Trained model for partitioning using informed random partitioning with tagged data of the Wiktionary and annotated corpora.
Can be run using text data, and POS tags, combining the output phrases for higher recall.
Achieves state-of-the-art performance with flexible application to any text-based corpora.

\cite{handler2016bag}
Bag of What? Simple Noun Phrase Extraction for Text Analysis
Uses patterns to match noun phrases using Finite State Transducer (FST) to find all matches for a given Finite State Grammar (FST), e.g., the pattern Adjective Noun Noun in the grammar.

\cite{mandera2015useful}
In \cite{van2016estimating} a review of the work by \cite{mandera2015useful} is provided: they evaluated how the performance of these computational approaches is influenced by the size of available norming data.
To that end, the 13,915 words in the \cite{warriner2013norms} norms were split into a training set and a test set, using different splits (e.g., 90/10 or 50/50 ).
Similarity indices between all words were obtained through applying LSA or PMI to a corpus comprising 385 million tokens.
These were then used to predict the valence, arousal, and dominance of words, with neighborhood parameter k set to 30 (the optimal value described by \cite{bestgen2012checking}).
They find that accuracy is somewhat reliant on the size of available norms. For example, when working with PMI-based similarity, increasing the training sample (i.e., the ratings that can contribute to the estimates) from 10 \% of the Warriner norms to 90 \% raises the correlation between the test sample and the norm ratings from .61 to .72 for valence, from .37 to .51 for arousal, and from .51 to .61 for dominance.
(They also investigated a number of other extrapolation methods, all of which showed a similar or lower accuracy.)

\cite{li2016understanding}
Understanding Neural Networks through Representation Erasure
\cite{hellrich2016bad}
Bad Company---Neighborhoods in Neural Embedding Spaces Considered Harmful
Random initialization undermines the reliability of the word2vec family of models for diachronic and synchronic tasks.
SVD$_\textrm{PPMI}$, as proposed by \cite{levy2014neural}, was highlighted to avoid these shortcomings but was not tested.

\cite{van2016estimating}
Estimating affective word covariates using word association data
Does just that for 14K dutch words, finding the best correlation between this method and human evaluation for $k$-NN algorithm (also tried ``Orientation towards Paradigm Words'').
For $k=10$ obtained correlations for valence, arousal, and dominance of .91, .84, and .85.
This performance is considerably better than was acheived by \cite{mandera2015useful} for English using corpus derived word similarity.

\cite{rothe2016ultradense}
Transform the embedding space via optimization of certain dimensions onto known semantic properties.
Awesome. Apply SGD to learn a transformation $Q$ that orthogonalizes the embedding matrix $A$ under the constraint of establishing a sentiment dimension.
